\section{Sprint 2: análisis de la solicitud}
\subsection{Propósito y objetivos}
El propósito fue procesar planes con texto o escaneados, calcular créditos y contenidos con las reglas vigentes y consultar precedentes para entregar una estimación explicable a revisión docente. El período planificado fue del 21 de octubre al 4 de noviembre de 2026, con revisión el 4 de noviembre. Las cinco historias fueron HU-07 a HU-11.

\subsection{Capacidad}
\begin{tabularx}{\textwidth}{|X|l|c|c|c|c|}\hline\rowcolor{grisTarjeta}\textbf{Integrante} & \textbf{Rol} & \textbf{Brutas} & \textbf{Coordinación} & \textbf{Técnicas} & \textbf{Planificado}\\\hline
Jair León & Scrum Master / Backend & 30 h & 1 h & 29 h & 46 h\\\hline Paul Acurio & Product Owner Proxy / IA & 30 h & 1 h & 29 h & 36 h\\\hline Alex Chisag & Frontend / UI-UX & 30 h & 1 h & 29 h & 32 h\\\hline Rolando Sebastián Miranda & Backend / QA-DevOps & 30 h & 1 h & 29 h & 39 h\\\hline Total & & 120 h & 4 h & 116 h & 153 h\\\hline
\end{tabularx}\captionof{table}{Capacidad del Sprint 2}
Las tareas sumaron 153 horas frente a una capacidad técnica de 116 horas. La planificación registró una diferencia de 37 horas como riesgo del sprint.

\subsection{Sprint Backlog}
\begin{tabularx}{\textwidth}{|c|X|c|c|}\hline\rowcolor{grisTarjeta}\textbf{Historia} & \textbf{Alcance} & \textbf{SP} & \textbf{Horas}\\\hline
HU-07 & Leer el plan analítico y confirmar datos dudosos & 8 & 36\\\hline HU-08 & Calcular con reglas y umbrales aprobados & 3 & 18\\\hline HU-09 & Mostrar proporción de créditos y alertas & 3 & 20\\\hline HU-10 & Leer planes escaneados y comparar contenidos & 8 & 40\\\hline HU-11 & Consultar casos anteriores y resoluciones & 5 & 39\\\hline Total & & 27 & 153\\\hline
\end{tabularx}\captionof{table}{Sprint Backlog del Sprint 2}

\subsection{Historias y pruebas de aceptación}
Las siguientes tarjetas condensan las historias incluidas en el informe entregado del Sprint 2. En la tabla de pruebas se resumen los escenarios de aceptación descritos en sus secciones.

\subsubsection{HU-07: Lectura del plan y confirmación}
\historia{HU-07}{Operador académico}{Leer el plan analítico y confirmar los datos dudosos}{Alta}{Media}{8 SP}{21 al 23 de octubre}{Jair León (5 h); Paul Acurio (8 h); Alex Chisag (13 h); Rolando Miranda (10 h)}{Como operador académico, quiero que el sistema lea créditos, horas y temas del PDF y me pida revisar lo que no leyó con seguridad, para evitar evaluar con datos equivocados.}{\item Dado un PDF con texto, cuando se envía, entonces se muestran los datos leídos y un aviso repetido no genera una evaluación duplicada.\item Dado un dato obligatorio ausente o confianza menor a 0,80, cuando termina la lectura, entonces queda pendiente de confirmación y no empieza el cálculo.\item Dado un PDF escaneado, cuando se procesa, entonces queda pendiente de lectura con IA y no se inventan datos.}{La lectura escaneada continúa en HU-10.}{M-01 y M-02}
\begin{longtable}{|p{1.45cm}|p{5.5cm}|p{1.9cm}|p{2.9cm}|p{0.95cm}|p{2.0cm}|}
\hline\rowcolor{grisTarjeta}\textbf{Número} & \textbf{Tarea} & \textbf{Tipo} & \textbf{Responsable} & \textbf{Horas} & \textbf{Fechas}\\\hline
\endhead
TA-07.1 & Recibir cada solicitud enviada una sola vez & Desarrollo & Jair León & 3 h & 21/10\\\hline
TA-07.2 & Preparar la IA local y las tablas del análisis & Otra & Sebastián Miranda & 5 h & 21/10\\\hline
TA-07.3 & Leer créditos, horas y unidades de un PDF con texto & Desarrollo & Paul Acurio & 5 h & 21/10\\\hline
TA-07.4 & Pantalla de bandeja de procesamiento & Desarrollo & Alex Chisag & 3 h & 21/10\\\hline
TA-07.5 & Servicios para consultar y confirmar los datos leídos & Desarrollo & Jair León & 2 h & 22/10\\\hline
TA-07.6 & Detectar PDF escaneados y lecturas poco seguras & Desarrollo & Paul Acurio & 3 h & 22/10\\\hline
TA-07.7 & Formulario para revisar y corregir los datos leídos & Desarrollo & Alex Chisag & 5 h & 22/10\\\hline
TA-07.8 & Pruebas unitarias de la lectura y del aviso repetido & Otra & Sebastián Miranda & 5 h & 22/10\\\hline
TA-07.9 & Conectar la confirmación con el Servicio de Análisis & Desarrollo & Alex Chisag & 5 h & 23/10\\\hline
\multicolumn{4}{|r|}{\textbf{Total}} & \textbf{36 h} & \\\hline
\caption{Tareas de la HU-07 (9 tareas)}
\end{longtable}

\begin{longtable}{|p{1.4cm}|p{3.7cm}|p{4.6cm}|p{5.1cm}|}
\hline\rowcolor{grisTarjeta}\textbf{N.º} & \textbf{Prueba} & \textbf{Entrada} & \textbf{Resultado esperado}\\\hline
\endhead
PA-07.1 & Lectura de un PDF con texto & Solicitud enviada con un plan analítico con texto & Se muestran los datos leídos con confianza 0,94 y el estado Datos leídos.\\\hline
PA-07.2 & Aviso de solicitud repetido & El mismo aviso llega dos veces & Se crea una sola evaluación.\\\hline
PA-07.3 & Lectura poco segura & PDF sin créditos legibles & Estado Confirmación pendiente; el cálculo no empieza.\\\hline
PA-07.4 & Confirmar con datos vacíos & Formulario con los créditos vacíos & Aviso Ingrese los créditos; no se guarda nada.\\\hline
PA-07.5 & PDF escaneado & PDF sin texto & Estado Pendiente de lectura con IA, sin datos inventados.\\\hline
\caption{Pruebas de aceptación de la HU-07 (definidas en el Sprint 2)}
\end{longtable}


\subsubsection{HU-08: Reglas y umbrales}
\historia{HU-08}{Docente evaluador}{Calcular con las reglas y umbrales aprobados}{Alta}{Media}{3 SP}{23 y 30 de octubre}{Jair León (10 h); Paul Acurio (3 h); Alex Chisag (2 h); Rolando Miranda (3 h)}{Como docente evaluador, quiero que el cálculo use las reglas vigentes aprobadas, para que el análisis sea consistente y explicable.}{\item Dada una solicitud confirmada, cuando se procesa, entonces el servicio consulta la versión vigente de las reglas.\item Dado un resultado calculado, cuando se presenta, entonces se identifican las reglas aplicadas.}{Consume las reglas del MS A.}{Reglas aplicadas en el análisis}
\begin{longtable}{|p{1.45cm}|p{5.5cm}|p{1.9cm}|p{2.9cm}|p{0.95cm}|p{2.0cm}|}
\hline\rowcolor{grisTarjeta}\textbf{Número} & \textbf{Tarea} & \textbf{Tipo} & \textbf{Responsable} & \textbf{Horas} & \textbf{Fechas}\\\hline
\endhead
TA-08.1 & Consultar las reglas vigentes al Servicio de Reglas & Desarrollo & Jair León & 10 h & 23/10\\\hline
TA-08.2 & Probar el cambio de un umbral y la caída del servicio & Mejora & Paul Acurio & 3 h & 23/10\\\hline
TA-08.3 & Pruebas de integración de la consulta de reglas & Otra & Sebastián Miranda & 3 h & 23/10\\\hline
TA-08.4 & Mostrar en las pantallas las reglas aplicadas & Desarrollo & Alex Chisag & 2 h & 30/10\\\hline
\multicolumn{4}{|r|}{\textbf{Total}} & \textbf{18 h} & \\\hline
\caption{Tareas de la HU-08 (4 tareas)}
\end{longtable}

\begin{longtable}{|p{1.4cm}|p{3.7cm}|p{4.6cm}|p{5.1cm}|}
\hline\rowcolor{grisTarjeta}\textbf{N.º} & \textbf{Prueba} & \textbf{Entrada} & \textbf{Resultado esperado}\\\hline
\endhead
PA-08.1 & Reglas en el resultado & Reglas en su versión 3; análisis de un expediente & El resultado muestra 80 ~\%, 15 ~\%, 0,75 y versión 3.\\\hline
PA-08.2 & Cambio de umbral & El administrador sube el parecido mínimo a 0,80 & Se usa 0,80 y la versión sube a 4.\\\hline
PA-08.3 & Servicio de Reglas apagado & Análisis de un expediente confirmado & Estado Pendiente de reglas con el error registrado.\\\hline
\caption{Pruebas de aceptación de la HU-08 (definidas en el Sprint 2)}
\end{longtable}


\subsubsection{HU-09: Créditos y alertas}
\historia{HU-09}{Docente evaluador}{Ver proporción de créditos y alertas}{Alta}{Media}{3 SP}{26 de octubre}{Jair León (7 h); Paul Acurio (5 h); Alex Chisag (3 h); Rolando Miranda (5 h)}{Como docente evaluador, quiero ver la proporción de créditos y las alertas antes de comparar, para valorar las condiciones de la solicitud.}{\item Dado un caso con créditos, cuando se analiza, entonces se muestra su proporción.\item Dado un caso fuera de las condiciones de vigencia o matrícula, cuando se analiza, entonces se presentan las alertas.}{Depende de los datos confirmados de HU-07 y las reglas de HU-08.}{Resumen de créditos y alertas}
\begin{longtable}{|p{1.45cm}|p{5.5cm}|p{1.9cm}|p{2.9cm}|p{0.95cm}|p{2.0cm}|}
\hline\rowcolor{grisTarjeta}\textbf{Número} & \textbf{Tarea} & \textbf{Tipo} & \textbf{Responsable} & \textbf{Horas} & \textbf{Fechas}\\\hline
\endhead
TA-09.1 & Calcular la proporción de créditos y los controles de vigencia y matrícula & Desarrollo & Jair León & 7 h & 26/10\\\hline
TA-09.2 & Pruebas de los controles con casos de la normativa & Mejora & Paul Acurio & 5 h & 26/10\\\hline
TA-09.3 & Panel de créditos y alertas & Desarrollo & Alex Chisag & 3 h & 26/10\\\hline
TA-09.4 & Pruebas con casos límite del cálculo de créditos & Otra & Sebastián Miranda & 5 h & 26/10\\\hline
\multicolumn{4}{|r|}{\textbf{Total}} & \textbf{20 h} & \\\hline
\caption{Tareas de la HU-09 (4 tareas)}
\end{longtable}

\begin{longtable}{|p{1.4cm}|p{3.7cm}|p{4.6cm}|p{5.1cm}|}
\hline\rowcolor{grisTarjeta}\textbf{N.º} & \textbf{Prueba} & \textbf{Entrada} & \textbf{Resultado esperado}\\\hline
\endhead
PA-09.1 & Créditos suficientes & 4 créditos de origen y 3 de destino & 100,00 ~\% con la etiqueta dato de apoyo.\\\hline
PA-09.2 & Créditos insuficientes & 2 créditos de origen y 3 de destino & 66,67 ~\%.\\\hline
PA-09.3 & Destino sin créditos & 4 frente a 0 & Mensaje de error; no se calcula.\\\hline
PA-09.4 & Materia de más de 10 años & Cambio de carrera; materia aprobada hace 12 años & Alerta de vigencia y paso a Validación de conocimientos.\\\hline
PA-09.5 & Tercera matrícula & Estudiante con tercera matrícula & Caso bloqueado con la alerta correspondiente.\\\hline
\caption{Pruebas de aceptación de la HU-09 (definidas en el Sprint 2)}
\end{longtable}


\subsubsection{HU-10: Lectura escaneada y comparación}
\historia{HU-10}{Docente evaluador}{Leer planes escaneados y comparar tema por tema}{Alta}{Alta}{8 SP}{27 al 29 de octubre}{Jair León (11 h); Paul Acurio (13 h); Alex Chisag (5 h); Rolando Miranda (11 h)}{Como docente evaluador, quiero que el sistema lea planes escaneados y compare sus temas, para revisar la cobertura de contenidos con evidencia.}{\item Dado un PDF escaneado, cuando se procesa, entonces se extrae información sin inventar los campos que no se reconocen.\item Dadas las materias de origen y destino, cuando se comparan, entonces se presenta el resultado por tema.}{Usa datos confirmados y reglas vigentes.}{Comparación de contenidos}
\begin{longtable}{|p{1.45cm}|p{5.5cm}|p{1.9cm}|p{2.9cm}|p{0.95cm}|p{2.0cm}|}
\hline\rowcolor{grisTarjeta}\textbf{Número} & \textbf{Tarea} & \textbf{Tipo} & \textbf{Responsable} & \textbf{Horas} & \textbf{Fechas}\\\hline
\endhead
TA-10.1 & Representar los temas con la IA local & Desarrollo & Jair León & 3 h & 27/10\\\hline
TA-10.2 & Leer PDF escaneados con OCR e IA sin inventar datos & Desarrollo & Paul Acurio & 4 h & 27/10\\\hline
TA-10.3 & Pantalla del estado de la lectura & Desarrollo & Alex Chisag & 2 h & 27/10\\\hline
TA-10.4 & Guardar los vectores de las materias de Software & Desarrollo & Sebastián Miranda & 3 h & 27/10\\\hline
TA-10.5 & Comparar temas y calcular el porcentaje de contenido & Desarrollo & Jair León & 4 h & 28/10\\\hline
TA-10.6 & Ajustar la comparación con casos reales & Mejora & Paul Acurio & 4 h & 28/10\\\hline
TA-10.7 & Tabla de cobertura por unidad & Desarrollo & Alex Chisag & 3 h & 28/10\\\hline
TA-10.8 & Pruebas unitarias del porcentaje de contenido & Otra & Sebastián Miranda & 5 h & 28/10\\\hline
TA-10.9 & Combinar dos materias de origen para un destino & Desarrollo & Jair León & 4 h & 29/10\\\hline
TA-10.10 & Pruebas funcionales de la combinación y de los temas faltantes & Mejora & Paul Acurio & 5 h & 29/10\\\hline
TA-10.11 & Pruebas de carga de la lectura con IA & Otra & Sebastián Miranda & 3 h & 29/10\\\hline
\multicolumn{4}{|r|}{\textbf{Total}} & \textbf{40 h} & \\\hline
\caption{Tareas de la HU-10 (11 tareas)}
\end{longtable}

\begin{longtable}{|p{1.4cm}|p{3.7cm}|p{4.6cm}|p{5.1cm}|}
\hline\rowcolor{grisTarjeta}\textbf{N.º} & \textbf{Prueba} & \textbf{Entrada} & \textbf{Resultado esperado}\\\hline
\endhead
PA-10.1 & PDF escaneado legible & PDF sin texto & Datos leídos; lo dudoso queda vacío para el operador.\\\hline
PA-10.2 & Cobertura de 6 de 7 unidades & 7 unidades de 12 h; 6 con parecido de 0,75 o más & 85,71 ~\% y la unidad 7 marcada como faltante.\\\hline
PA-10.3 & Parecido bajo el umbral & Unidad con parecido 0,74 & La unidad queda como faltante.\\\hline
PA-10.4 & Dos materias para un destino & Física I y Física II para Física & Los temas se suman sin contarse dos veces.\\\hline
PA-10.5 & Misma materia en dos destinos & Física I usada en dos destinos & El sistema lo rechaza.\\\hline
\caption{Pruebas de aceptación de la HU-10 (definidas en el Sprint 2)}
\end{longtable}


\subsubsection{HU-11: Precedentes}
\historia{HU-11}{Docente evaluador}{Consultar casos anteriores y sus resoluciones}{Media}{Media}{5 SP}{29 de octubre al 3 de noviembre}{Jair León (13 h); Paul Acurio (7 h); Alex Chisag (9 h); Rolando Miranda (10 h)}{Como docente evaluador, quiero consultar casos anteriores idénticos y su resolución, para disponer de antecedentes durante la revisión.}{\item Dada una solicitud con materias identificables, cuando se consulta, entonces se muestran precedentes coincidentes y su resolución.\item Dado que no hay precedentes, cuando se consulta, entonces se informa que no existen coincidencias.}{La búsqueda de casos anteriores aporta contexto para la revisión.}{Consulta de precedentes}
\begin{longtable}{|p{1.45cm}|p{5.5cm}|p{1.9cm}|p{2.9cm}|p{0.95cm}|p{2.0cm}|}
\hline\rowcolor{grisTarjeta}\textbf{Número} & \textbf{Tarea} & \textbf{Tipo} & \textbf{Responsable} & \textbf{Horas} & \textbf{Fechas}\\\hline
\endhead
TA-11.1 & Pantalla de consulta de precedentes & Desarrollo & Alex Chisag & 5 h & 29/10\\\hline
TA-11.2 & Guardar precedentes y buscarlos por coincidencia exacta & Desarrollo & Jair León & 8 h & 30/10\\\hline
TA-11.3 & Cargar resoluciones anteriores de prueba & Mejora & Paul Acurio & 3 h & 30/10\\\hline
TA-11.4 & Pruebas de la búsqueda de precedentes & Otra & Sebastián Miranda & 5 h & 30/10\\\hline
TA-11.5 & Unir todos los pasos del análisis & Desarrollo & Jair León & 3 h & 02/11\\\hline
TA-11.6 & Verificar que un precedente no salte los controles & Mejora & Paul Acurio & 2 h & 02/11\\\hline
TA-11.7 & Mensajes de error claros en la interfaz & Corrección & Alex Chisag & 2 h & 02/11\\\hline
TA-11.8 & Prueba completa del flujo entre servicios & Otra & Sebastián Miranda & 3 h & 02/11\\\hline
TA-11.9 & Ordenar el código y cerrar la parte técnica & Mejora & Jair León & 2 h & 03/11\\\hline
TA-11.10 & Repetir las pruebas funcionales y guardar evidencias & Mejora & Paul Acurio & 2 h & 03/11\\\hline
TA-11.11 & Revisar el aspecto de las pantallas y preparar la demostración & Mejora & Alex Chisag & 2 h & 03/11\\\hline
TA-11.12 & Agregar las pruebas a la integración continua & Otra & Sebastián Miranda & 2 h & 03/11\\\hline
\multicolumn{4}{|r|}{\textbf{Total}} & \textbf{39 h} & \\\hline
\caption{Tareas de la HU-11 (12 tareas)}
\end{longtable}

\begin{longtable}{|p{1.4cm}|p{3.7cm}|p{4.6cm}|p{5.1cm}|}
\hline\rowcolor{grisTarjeta}\textbf{N.º} & \textbf{Prueba} & \textbf{Entrada} & \textbf{Resultado esperado}\\\hline
\endhead
PA-11.1 & Coincidencia exacta & Física I y Física II para Física, misma universidad & Se muestran la resolución, las materias y la universidad.\\\hline
PA-11.2 & Sin coincidencia & Bases de Datos I, Universidad de Cuenca & Aviso Sin precedentes registrados.\\\hline
PA-11.3 & Caso solo parecido & Física I sola para Física & No se muestra como coincidencia.\\\hline
PA-11.4 & Precedente y más de 10 años & Precedente favorable; materia de hace 12 años & Igual pide validación de conocimientos y revisión docente.\\\hline
PA-11.5 & Falla la búsqueda & Servicio de precedentes sin respuesta & Aviso; los cálculos de créditos y temas continúan.\\\hline
\caption{Pruebas de aceptación de la HU-11 (definidas en el Sprint 2)}
\end{longtable}


\subsection{Resultado del Sprint 2}
El informe del Sprint 2 registra la demostración del incremento con lectura, confirmación de datos, reglas, créditos, contenidos y precedentes. Una prueba de aceptación de HU-11 (PA-11.5) quedó pendiente y otra quedó parcial según la evidencia de ejecución. En el trabajo posterior se abordó la falla del servicio y el reintento automático, descritos en la sección de correcciones.
